% Figure: membrane blood oxygenator. Left panel sketches a hollow fibre with
% oxygen flowing through the lumen and blood over the shell side. Right panel
% plots the oxygen partial pressure across the gas film, the membrane wall,
% and the blood film, with the steep drop on the blood side marking the
% controlling resistance.
\begin{figure}[htbp]
\centering
\begin{tikzpicture}[>=stealth, font=\small]
% left panel: fibre cross-section schematic
\begin{scope}[shift={(0,0)}]
  % outer blood region
  \fill[engred!12] (-1.6,-1.6) rectangle (1.6,1.6);
  % membrane wall (annulus drawn as two circles)
  \fill[enggray!35] (0,0) circle (0.95);
  \fill[white] (0,0) circle (0.7);
  % oxygen lumen
  \fill[engblue!18] (0,0) circle (0.7);
  \draw[thick] (0,0) circle (0.95);
  \draw[thick] (0,0) circle (0.7);
  \node[engblue] at (0,0) {O$_2$};
  \node[enggray, font=\scriptsize, anchor=south] at (0,0.95) {membrane};
  \node[engred, font=\scriptsize] at (1.15,1.25) {blood};
  % arrows: oxygen out of lumen into blood
  \draw[->, engblue, thick] (0.7,0) -- (1.5,0);
  \draw[->, engblue, thick] (-0.7,0) -- (-1.5,0);
  \node[font=\scriptsize, anchor=south] at (-2.0,-2.0) {fibre cross-section};
\end{scope}
% right panel: partial-pressure profile across the three films
\begin{scope}[shift={(6.4,0)}]
  \draw[->] (-2.2,-1.6) -- (2.4,-1.6) node[right, font=\scriptsize] {position};
  \draw[->] (-2.2,-1.6) -- (-2.2,1.6) node[above, font=\scriptsize] {$p_{\mathrm{O_2}}$};
  % zone dividers: gas film | membrane | blood film
  \draw[enggray, dashed] (-1.2,-1.6) -- (-1.2,1.4);
  \draw[enggray, dashed] (-0.4,-1.6) -- (-0.4,1.4);
  % gas side high, small drop across gas film and membrane, big drop in blood
  \draw[engblue, very thick]
    (-2.2,1.2) -- (-1.2,1.1) -- (-0.4,0.95) -- (2.2,-0.9);
  % labels
  \node[font=\scriptsize, anchor=south] at (-1.7,1.15) {gas};
  \node[font=\scriptsize, anchor=south] at (-0.8,1.05) {wall};
  \node[font=\scriptsize, anchor=south, engred] at (1.0,-0.2) {blood film};
  \node[font=\scriptsize, anchor=west] at (-2.15,1.35) {pure O$_2$};
  \node[font=\scriptsize, anchor=west] at (1.4,-1.0) {blood O$_2$};
\end{scope}
\end{tikzpicture}
\caption[Membrane blood oxygenator: fibre and oxygen profile]{A membrane
blood oxygenator. Oxygen flows through the lumen of a microporous hollow
fibre while blood passes over the shell side; oxygen crosses the membrane
wall into the blood and carbon dioxide crosses the other way. The right
panel plots the oxygen partial pressure across the three resistances in
series: the gas film and the membrane wall drop little, because oxygen
moves fast through gas and through the open pores, while the blood film
carries almost the whole partial-pressure drop, because oxygen diffuses
slowly through plasma and is taken up slowly by the red cells. The blood
film is the controlling resistance, which is why a commercial unit carries
several times the bare-estimate fibre area as margin against blood
maldistribution and against a plasma layer building on the fibres in a
long run.}
\label{fig:vol04ch07:oxygenator}
\end{figure}
